\documentclass[11pt,a4paper]{article}
\usepackage[T1]{fontenc}
\usepackage{lmodern,amsmath,amssymb,amsthm,booktabs,array}
\usepackage[margin=27mm]{geometry}
\usepackage[hidelinks]{hyperref}
\usepackage{microtype}
\hypersetup{pdftitle={Matrix gradients and critical observations},
pdfauthor={Lluis Eriksson},pdfsubject={Consolidated draft appendix to the coarse information manuscript},
pdfkeywords={finite gradients, inverse moments, spectral repair, effective resistance}}
\newtheorem{theorem}{Theorem}
\newtheorem{proposition}[theorem]{Proposition}
\newtheorem{example}[theorem]{Example}
\newcommand{\R}{\mathbb R}
\newcommand{\tr}{\operatorname{tr}}
\newcommand{\rank}{\operatorname{rank}}
\newcommand{\im}{\operatorname{im}}
\newcommand{\id}{\mathrm I}
\newcommand{\Nlow}{N_{<c}}
\DeclareMathOperator{\diag}{diag}
\title{Matrix gradients and critical observations\\
\large Consolidated draft on signed channels and rectangular means}
\author{Lluis Eriksson}
\date{7 November 2026}
\begin{document}
\maketitle
\begin{abstract}
We consolidate finite matrix gradient estimates, spectral repair bounds and
signed channel examples into one supporting appendix. An operator-valued
Cwikel estimate gives an inverse moment bound with explicit counting-measure
normalization. Finite-rank repair produces a capped moment criterion that
admits zero and negative channels. A separate effective-resistance argument
constructs stable signed families with uniform observation rank and penalty,
and a small Schur matrix detects coupled instabilities. A critical rectangular mean obstruction is also transferred to actual flat coexact cochains. An explicit torus
example shows that individually repairable channels need not remain
repairable together by a scalar differentiable observation. The underlying
spectral and resistance methods are classical. The statements concern the
defined finite gradient form; they establish neither an interacting
Yang--Mills form identity nor a continuum mass gap.
\end{abstract}

\paragraph{Status and provenance.}
This is a grouped appendix candidate for the existing integrated manuscript,
not a separate deposit. Its analytical sources and finite evidence are pinned
at repository commit \texttt{1e3ede17}; full identifiers and execution records
are in Section~\ref{sec:evidence}. The published working paper v6 and its
assessment report 154 remain separate artifacts. This new PDF has no inherited
assessment score, independent review or Lean certificate.
The parent is
\href{https://airr.science/papers/ARR-2026-7EPM3V3DP796BBRR/versions/v6/}
{AIRR working paper ARR-2026-7EPM3V3DP796BBRR version 6}; its unchanged PDF
SHA256 is
\begin{center}\scriptsize
\texttt{ccf53e0afc9819c3fc9b2008de9076a60f60d23ef6ac2c4ccd82adca30cbf389}.
\end{center}

\section{Actual finite objects and observation scope}
Let $D\geq3$, $n\geq2$, $h\geq1$, $V=n^D$, $r=Dh$, and $p=D/2$.
On $X=(\mathbb Z/n\mathbb Z)^D$ use the counting inner product. For
$u:X\to\R^h$, define
\[
 du(x)=(u(x+e_i)-u(x))_{i=1}^D,\qquad
 Q_A(u)=\sum_x du(x)^T A(x)du(x),\quad L_A=d^*A d,
\]
where $A(x)$ is a real symmetric $r\times r$ matrix. A channel is one local
coefficient eigenvector, counted with multiplicity. Distinct matrices need
not commute. All derivatives have unit spacing. A physical derivative $nd$
requires an explicit rescaling of every energy and threshold.

The number $\Nlow(L_A)$ counts all eigenvalues strictly below $c>0$,
including negative and zero ones, with ordinary real multiplicity. Spectral
coordinates on this low space give a real observation $C$ and penalty
\begin{equation}\label{eq:penalty}
 \Phi(y)=\sum_{\lambda<c}(c-\lambda)y_\lambda^2,
 \qquad Q_A(u)+\Phi(Cu)\geq c\|u\|^2.
\end{equation}
This is an existence result for an actual, generally nonlocal observation.
It does not construct the unresolved local rectangular observation.
For signed forms, $c\|y\|^2$ alone need not be sufficient.

\section{The positive matrix estimate}
Put
\[
 \alpha_D=\frac D2\left(\frac D{D-2}\right)^{D-1}
                 \left(\frac34\right)^D.
\]
\begin{theorem}\label{thm:positive}
If every $A(x)$ is positive definite, then
\begin{equation}\label{eq:positive}
 \Nlow(L_A)\leq h+\alpha_D c^p\sum_x\tr A(x)^{-p}.
\end{equation}
\end{theorem}
\begin{proof}
On the mean-zero space $H_0$, write $L_0=d^*d$ and $J=dL_0^{-1}$.
For any $b\in H_0$, the flow $j_*=A dL_A^{-1}b$ has $d^*j_*=b$.
Every other feasible flow $j=j_*+k$ has $d^*k=0$, and expansion gives
\[
 \langle j,A^{-1}j\rangle
 =\langle b,L_A^{-1}b\rangle+\langle k,A^{-1}k\rangle.
\]
Taking the actual feasible flow $Jb$ proves
$L_A^{-1}\leq J^*A^{-1}J$. No coefficients commute in this argument.

The unitary counting Fourier transform on $r$ components has
$L_1(\ell_2^r)\to L_\infty(\ell_2^r)$ norm $V^{-1/2}$. The symbol of $J$
is $z(k)/\lambda(k)$ with
$z_i(k)=e^{2\pi i k_i/n}-1$ and
$\lambda(k)=4\sum_i\sin^2(\pi k_i/n)$; set $J(0)=0$.
It is rectangular $r\times h$. Embed $\R^h$ isometrically in $\R^r$ by $I$
and use the square symbol $G(k)=J(k)I^*$. Thus $\|G(k)\|=\lambda(k)^{-1/2}$
for $k\ne0$. Writing $\mathcal F_r$ for the counting transform, the actual
factorization is $J=\mathcal F_r^*M_G\mathcal F_r I$.
Centered frequency representatives $b$ satisfy
$\lambda(k)\geq16|b|^2/n^2$. Counting the cube containing
$0<|b|<n/(4t)$ gives
\[
 \sup_{t>0}t^D\#\{k:\|G(k)\|>t\}\leq(3/4)^D V.
\]
When the radius is at most one the set is empty; otherwise
$(2\lfloor R\rfloor+1)^D\leq(3R)^D$ suffices.

Apply Frank's operator-valued Cwikel theorem~\cite[Theorem 3.1]{Frank}
with $q=D$, $f=A^{-1/2}$, this square multiplier and counting measures.
Its constant is $\frac D2(\frac D{D-2})^{D-1}V^{-1}$, so the displayed
frequency $V$ cancels exactly once. It bounds the number of singular values
of $A^{-1/2}J$ strictly above $s$ by
$\alpha_Ds^{-D}\sum_x\tr A(x)^{-p}$. The right factor $\mathcal F_r I$ is an isometry
and cannot increase singular-value counts. Apply min-max to the inverse
comparison at threshold $1/c$ on $H_0$, then add the $h$ constant modes.
Complexification of a real symmetric matrix preserves multiplicities.
\end{proof}

\paragraph{Input and interpretation.}
The operator-valued estimate is classical; this proof fixes its finite
normalization, padding and actual form dictionary. It is not a new general
Cwikel theorem. A matrix arithmetic mean alone does not imply its inverse
moment hypothesis. For the formal flat cochain application one may set
$h=D(N_c^2-1)$, but an interacting Hessian or weighted curl/divergence form
must first be proved equal or comparable to this $Q_A$.

\section{Spectral repair and the capped budget}
\begin{theorem}\label{thm:repair}
Let $A(x)$ now be arbitrary real symmetric matrices. In local orthonormal
eigenbases choose a set $S$ of channels containing all nonpositive ones.
Let $P_S$ be their orthogonal projection and $R(S)=\rank(P_Sd)$. Then
\begin{equation}\label{eq:repair}
 \Nlow(L_A)\leq\min\left\{hV,\ h+R(S)+
          \alpha_Dc^p\sum_{j\notin S}\lambda_j^{-p}\right\}.
\end{equation}
Writing $\nu_0$ for the total multiplicity of nonpositive local eigenvalues,
one consequently has
\begin{equation}\label{eq:capped}
 \Nlow(L_A)\leq\min\{hV,\ h+F_c(A)\},\quad
 F_c(A)=\nu_0+\sum_{\lambda_j>0}
                  \min\{1,\alpha_Dc^p\lambda_j^{-p}\}.
\end{equation}
\end{theorem}
\begin{proof}
Replace selected eigenvalues by $T$ larger than zero and all of them.
The resulting field $A_T$ is positive definite and
\[
 L_{A_T}-L_A=d^*P_S(T\id-A)P_Sd\geq0,
 \quad\rank(L_{A_T}-L_A)=R(S).
\]
The rank equality follows because $T\id-A$ is strictly positive on
$\im P_S$. Finite-rank min-max gives
$\Nlow(L_A)\leq\Nlow(L_{A_T})+R(S)$: intersect the low space of $L_A$
with the kernel of the update and the orthogonal complement of the low
space of $L_{A_T}$. A dimension larger than $R(S)$ would give a nonzero
vector whose identical quadratic forms are both strictly below and at least
$c\|u\|^2$. Theorem~\ref{thm:positive} and $T\to\infty$ prove
\eqref{eq:repair}. For \eqref{eq:capped}, use $R(S)\leq|S|$ and select
all positive channels with $\alpha_Dc^p\lambda_j^{-p}\geq1$ as well.
\end{proof}

\paragraph{A strict extension of the full moment criterion.}
Let $m\geq1$ and $T_m=\max\{1,V/(1+m)^p\}$. If $F_c(A)\leq KT_m$,
then $\Nlow(L_A)\leq(h+K)T_m$. Every former SPD budget
\[
 \frac1{rV}\sum_x\tr((mA(x)^{-1})^p)\leq M
\]
implies this capped budget, with $K=2^p\alpha_Dc^prM$.
Conversely choose $m=n^2$, one actual diagonal edge channel
$\epsilon=n^{-4}$, and all others $m$. Now $T_m=1$ and
$F_c(A)\leq1+\alpha_Dc^pr$, whereas the normalized full inverse moment is
at least $n^{2D}/r$ and diverges. These bounds are derived for this explicit
family; no new moment bound is assumed.

\paragraph{Sparse contamination and its limitation.}
Selecting every channel with $\lambda<m$ gives
\[
 \Nlow(L_A)\leq h+R_m+\alpha_Dc^prVm^{-p},
 \qquad R_m=\rank(\boldsymbol1_{(-\infty,m)}(A)d).
\]
An actual selected-channel count $k_m\leq KT_m$ suffices since $R_m\leq k_m$.
If $A(x)\geq m\id$ outside a specified set $B$, then $k_m\leq r|B|$.
Neither condition follows from an arithmetic mean. This particular
prescribed-threshold count criterion does not subsume every full moment
family: $A=(m/2)\id$ has normalized inverse moment $2^p$ but every channel
lies below $m$. The two particular sufficient criteria are incomparable;
the general capped criterion \eqref{eq:capped} does include the old one.

\paragraph{Actual rank and unavoidable costs.}
For diagonal direction/component channels the selected rows are graph
incidences. If $\kappa_a$ counts all components, including isolated vertices,
of the selected-edge graph for component $a$, then
$R(S)=\sum_a(V-\kappa_a)$: its kernel consists precisely of fields constant
on those components. Parallel rows at $n=2$ do not change that kernel.
This formula is not applied to rotated local eigenchannels.
Isolating $q$ nonadjacent vertices by zeroing their incident edges gives
at least $q+1$ kernel modes. Replacing these zeros by
$0<\epsilon<c/(2D)$ yields a $q$-dimensional strict low subspace of delta
functions. The defect-count cost cannot simply be deleted.

In \eqref{eq:penalty}, a generic bound is $K\leq c+4Db$ with
$b=\max_x\|A_-(x)\|$, since $\|du\|^2\leq4D\|u\|^2$.
One negative channel $-B$ with all other channels $m$, at $n\geq3,h=1$,
has the trial vector $v=\delta_y-\delta_x$ on that edge and
\[
 \frac{Q_A(v)}{\|v\|^2}=-2B+(2D-1)m.
\]
Any normalized linear observation with penalty bounded by $K\|y\|^2$
therefore needs $K\geq c+2B-(2D-1)m$. Sparsity alone does not bound penalty.

\section{Stable signed families from actual resistance}
This section is independent of the Cwikel input. Its mechanism is classical
effective resistance and finite-rank congruence~\cite{Zelazo,Spielman}.
Let $e$ be one actual direction/component row and $v=d^*e$. With all
coefficients $m$ except this one set to $-B$, $B\geq0$,
\[
 H=mL_0-(m+B)vv^*,\qquad L_0=d^*d.
\]
On $H_0$ let $G=L_0^{-1}$. Translations and coordinate permutations make
every scalar row leverage equal. Their sum is
$\tr(dGd^*)=V-1$, so
\begin{equation}\label{eq:rho}
 \rho=v^*Gv=\frac{V-1}{DV}.
\end{equation}
Components decouple; the same value holds for every $h$. The $DV$ rows
are counted individually even at $n=2$.

\begin{proposition}\label{prop:phase}
The exact threshold is $B_*=m(DV/(V-1)-1)$. On $H_0$, $H$ is positive
definite for $B<B_*$, has one zero direction at equality, and has exactly
one negative direction above it. If $m=n^2$, $B=\theta m$ and
$0\leq\theta<D-1$, its mean-zero gap is at least
\[
 g=16(D-1-\theta)/D>0.
\]
For every $0<c\leq g$, the normalized mean observation
$C(u)=V^{-1/2}\sum_xu(x)\in\R^h$ and $\Phi(y)=c\|y\|^2$ give coercivity
with optimal rank $h$, uniformly in $n$.
\end{proposition}
\begin{proof}
Congruence by $G^{1/2}$ gives $m\id-(m+B)ww^*$ with $\|w\|^2=\rho$,
proving the phase claim. The equality witness is $Gv$.
Cauchy--Schwarz in the $L_0$ energy gives
$|v^*u|^2\leq\rho\langle u,L_0u\rangle$, hence
$H\geq[m-(m+B)\rho]L_0$ on $H_0$; equality holds on $Gv$.
For the stated scaling the coefficient is at least
$m(D-1-\theta)/D$. The actual first nonconstant Fourier eigenvalue of $L_0$
is $4\sin^2(\pi/n)\geq16/n^2$. This proves the gap. Decompose into
constants and $H_0$ to prove the displayed observation inequality.
The constant kernel forces rank at least $h$.
\end{proof}
For $\theta>0$, negative coefficient depth grows as $\theta n^2$, yet the
penalty is uniformly $c$. A bounded-depth hypothesis is sufficient for the
generic estimate, not necessary for every stable signed family. The actual
form is positive despite an indefinite coefficient field; this does not
contradict the fixed-$m$, $B\to\infty$ penalty obstruction above.

\section{Coupled channels and a nonlinear observation obstruction}
Choose $q$ actual rows of $d$ as $Z$, and $W=\diag(w_j)>0$. The modified
form is exactly $H=mL_0-Z^*WZ$. Define
\[
 \Gamma=ZGZ^*,\qquad \mathcal C=mW^{-1}-\Gamma,\quad N_0=h(V-1).
\]
Schur elimination of either positive diagonal block in
$\left[\begin{smallmatrix}mL_0&Z^*\\Z&W^{-1}\end{smallmatrix}\right]$
on $H_0\oplus\R^q$ proves
\begin{equation}\label{eq:inertia}
 n_-(H)=n_-(\mathcal C),\quad n_0(H)=h+n_0(\mathcal C),\quad
 n_+(H)=N_0-q+n_+(\mathcal C).
\end{equation}
Indeed the two congruences give the identity
\[
 \operatorname{inertia}(H|_{H_0})+(q,0,0)
 = (N_0,0,0)+\operatorname{inertia}(\mathcal C/m).
\]
No disjoint-cycle or independent-row assumption occurs. In particular,
$H$ is positive semidefinite with exactly $h$ constant zeros if and only if
$\mathcal C$ is positive definite. General signed matrix criteria preexist
in~\cite{Chen}; our $\Gamma$ uses the positive reference inverse.

\begin{example}\label{ex:coupled}
Take $D=3,n=2,h=1,m=4$, and change the first two direction rows at one
vertex separately or together to coefficient $-84/11$. Then
\[
 \Gamma=\begin{pmatrix}7/24&5/48\\5/48&7/24\end{pmatrix},
 \quad w=128/11,\quad w/m=32/11.
\]
The diagonal follows from \eqref{eq:rho}; the off-diagonal Walsh sum is
$(1/8)(1/2+1/3)=5/48$, with both forward orientations counted.
Each single channel is below its critical magnitude $68/7$ and has full
inertia $(7,0,1)$. Together the reduced matrix has diagonal $5/96$ and
off-diagonal $-5/48$, so its eigenvalues are $-5/96,5/32$.
Equation~\eqref{eq:inertia} gives joint inertia $(6,1,1)$.

Each single form has mean-zero gap at least
$16(1-(32/11)(7/24))=80/33>1$. The same scalar mean observation with
$\Phi(y)=|y|^2$ repairs either single form at $c=1$. The joint form cannot
be repaired by any scalar observation $F:\R^8\to\R$ differentiable at zero
with $F(0)=0$ and $0\leq\Phi(z)\leq K|z|^2$ near zero. Its negative
eigenvector and constant vector span a two-dimensional nonpositive space.
The derivative of $F$ has a nonzero kernel vector $v$ there. Then
$F(tv)=o(t)$ and $\Phi(F(tv))=o(t^2)$, whereas $Q_A(tv)\leq0$, contradicting
coercivity at $c=1$ for small nonzero $t$.
\end{example}
This refutes the shortcut from individual to joint control, not the full
matrix resistance criterion or any theorem retaining a disjoint-cycle
hypothesis. The $q$-row perturbation also gives
$\Nlow(H)\leq h+q$ for $0<c\leq16m/n^2$. Extra zero or negative modes
provide the necessary rank lower bound in \eqref{eq:inertia}; they do not
alone bound the penalty coefficient.

\section{Critical constants for rectangular mean observations}
For a disjoint partition $\mathcal P$ of a finite counting space, its
normalized means are $(Cf)_R=|R|^{-1/2}\sum_{x\in R}f(x)$. They satisfy
$CC^*=\id$ and $C^*C=P$, the orthogonal projection onto functions constant
on each cell. This concrete normalization matters for the penalty.

\begin{theorem}\label{thm:critical}
Let $L=L^*$ on a finite real Hilbert space, $c>0$, and $C$ a coisometry
onto $\R^q$, $q\geq1$. On $\mathcal K=\ker C$ and
$\mathcal O=\im C^*$, write
\[
 L-c\id=\begin{pmatrix}A_0&X\\X^*&B_0\end{pmatrix}.
\]
A finite $k\geq0$ satisfying $L+kC^*C\geq c\id$ exists if and only if
\begin{equation}\label{eq:compatibility}
 A_0\geq0,\qquad X^*(\ker A_0)=0.
\end{equation}
When these hold, its smallest value is
\[
 k_{\min}=\max\{0,\lambda_{\max}(X^*A_0^\dagger X-B_0)\}.
\]
Here $A_0^\dagger$ is the Moore--Penrose inverse on $\mathcal K$.
\end{theorem}
\begin{proof}
The restriction to $\mathcal K$ forces $A_0\geq0$. For $z\in\ker A_0$
the augmented form at $tz+y$ has linear term $2t\langle z,Xy\rangle$
and no quadratic term in $t$. Nonnegativity for every real $t$ forces
$X^*z=0$. This is equivalent to $\im X\subseteq\im A_0$. Under that
condition the full form at $z+y$ is exactly
\[
 \langle z+A_0^\dagger Xy,A_0(z+A_0^\dagger Xy)\rangle
 +\langle y,(B_0+k\id-X^*A_0^\dagger X)y\rangle.
\]
Finite spectral decomposition proves the criterion and the minimum.
\end{proof}
This is the classical generalized Schur criterion~\cite[Theorem 1]{Higham},
applied to actual means; it is not a new general matrix theorem. Strict
positivity of $A_0$ suffices, whereas a weak bound at the target requires
the extra compatibility in~\eqref{eq:compatibility}.

\begin{example}\label{ex:rectangle}
On $(\mathbb Z/2\mathbb Z)^3$ take the same scalar coefficient in every
gradient direction,
\[
 a(x)=6+(-1)^{x_1}+4(-1)^{x_2},\qquad L=d^*ad.
\]
Its values are $1,3,9,11$, each twice, and its mean is $m=6$; hence
$1\leq a\leq2m$. Choose the four actual rectangles
$R_{yz}=\{0,1\}\times\{y\}\times\{z\}$ and
$(Cf)_{yz}=(f(0,y,z)+f(1,y,z))/\sqrt2$.
The best coercivity constant of $L$ on $\ker C$ is $16$, but no finite
quadratic penalty for these means attains the same constant on the full
space. Every $0<c<16$ is attainable. Thus the supremum $16$ is not attained.
\end{example}
\begin{proof}
For $S\subseteq\{1,2,3\}$ let
$\chi_S(x)=(-1)^{\sum_{i\in S}x_i}/\sqrt8$. These are an orthonormal real
Walsh basis. Retain all $24$ oriented forward rows, including opposite rows
at period two. The actual adjoint calculation gives
\begin{align*}
 L\chi_S={}&24|S|\chi_S
 +4\bigl(|S|-\boldsymbol1_{1\in S}\bigr)\chi_{S\triangle\{1\}}\\
 &+16\bigl(|S|-\boldsymbol1_{2\in S}\bigr)\chi_{S\triangle\{2\}}.
\end{align*}
Indeed $d_i\chi_S=-2\chi_S$ when $i\in S$, and the corresponding
adjoint term is $2(a+T_i a)\chi_S$; absent directions contribute zero.
In the kernel basis $(\chi_1,\chi_{12},\chi_{13},\chi_{123})$ the compression
has the two blocks
\[
 \begin{pmatrix}24&16\\16&48\end{pmatrix},\qquad
 \begin{pmatrix}48&32\\32&72\end{pmatrix}.
\]
After subtracting $16\id$, the first is positive semidefinite with kernel
vector $(-2,1)$, and the second is positive definite, with first pivot $32$
and determinant $768$. This proves the complete kernel bound, not just a
bound on selected test vectors.

Set $v=-2\chi_1+\chi_{12}$ and $u=\chi_2$. Direct evaluation gives
\[
 Cv=0,\quad\|v\|^2=5,\quad\|Cu\|^2=\|u\|^2=1,
 \quad(L-16\id)v=4u,\quad\langle u,(L-16\id)u\rangle=8.
\]
The augmented form on $u+tv$ is therefore $8+8t+k$, negative for
$t<-(8+k)/8$. No finite $k$ works. For $c<16$ the compressed block is
strictly positive and Theorem~\ref{thm:critical} supplies a finite penalty;
for $c>16$, $v$ violates coercivity already in the kernel.
\end{proof}

\paragraph{A quantitative penalty obstruction.}
Write $\delta=16-c>0$. On the orthonormal pair $(v/\sqrt5,u)$ the
augmented form has matrix
\[
 \begin{pmatrix}\delta&4/\sqrt5\\4/\sqrt5&8+\delta+k\end{pmatrix}.
\]
Its determinant gives the necessary bound
\begin{equation}\label{eq:blowup}
 k\geq\max\{0,16/(5\delta)-8-\delta\}.
\end{equation}
This lower bound need not equal the full optimum. It proves divergence
as the target approaches $16$ from below without numerical extrapolation.
If $F(0)=0$, $DF(0)=C$ and $0\leq\Phi(y)\leq k\|y\|^2$ near zero,
local coercivity would imply $L+kC^*C\geq c\id$ by substituting $sf$,
dividing by $s^2$ and taking $s\to0$. Higher order terms with this same
derivative do not repair the failure at $c=16$.

\section{Transfer to actual flat coexact cochains}
The preceding obstruction also holds on the actual finite coexact domain,
with a proved form dictionary. Fix $N_c\geq2$ and an actual coordinate unit
$e\in\R^{N_c^2-1}$. At the trivial background, a physical one-cochain is
$A_j(x)\in\R^{N_c^2-1}$, $j=1,2,3$, with counting inner product. Use one
unordered plaquette copy for each $i<j$:
\[
 (d_1 A)_{ij}=d_iA_j-d_jA_i,\qquad
 \operatorname{div}A=\sum_j(A_j-T_j^{-1}A_j).
\]
Expansion of curl squares and cyclic summation by parts cancel the cross
terms against divergence, proving the all-cochain identity
\[
 \|d_1A\|^2+\|\operatorname{div}A\|^2=\sum_{i,j,x}\|d_iA_j(x)\|^2.
\]
These are the actual repository conventions, not a supplied comparison.
Take the actual componentwise repair coefficient $w(i,(x,j))=a(x)-1\geq0$.
The total Hodge-plus-repair form is consequently
\begin{equation}\label{eq:physical-dictionary}
 \begin{aligned} Q_w(A)={}&\|d_1A\|^2+\|\operatorname{div}A\|^2+
 \sum_{i,j,x}w(x)\|d_iA_j(x)\|^2
 \\ ={}&\sum_{i,j,x}a(x)\|d_iA_j(x)\|^2. \end{aligned}
\end{equation}
It is a direct sum of $3(N_c^2-1)$ scalar copies. Componentwise normalized
means on the same four bars have output dimension $12(N_c^2-1)$, and
the whole ambient kernel has coercivity constant $16$ by the two blocks
in Example~\ref{ex:rectangle}. Restriction to coexact fields preserves the
inequality; coisometry on the restricted domain is not assumed.

\begin{proposition}\label{prop:coexact}
At target $16$, no finite quadratic penalty for these actual componentwise
rectangle means supplies coercivity even on $\im d_1^*$. The same is true
for local nonlinear observations with derivative equal to those restricted
means and a locally quadratic penalty. On that domain every $0<c<16$ is
attainable, and the necessary bound~\eqref{eq:blowup} still applies.
\end{proposition}
\begin{proof}
Construct actual fields $U_3=ue$, $V_3=ve$ from Example~\ref{ex:rectangle},
with all other components zero. They are isometric copies of $u,v$ and
independent of $x_3$, so their actual backward divergence vanishes.
More explicitly, use unordered two-cochains with only
\[
 B^U_{23}=-\tfrac12\chi_2 e,\qquad
 B^V_{13}=(\chi_1-\tfrac12\chi_{12})e
\]
nonzero. At side two $d_i^*=T_i^{-1}-\id=T_i-\id$. The adjoint of the
displayed curl puts $d_i^*B_{ij}$ in component $j$ and $-d_j^*B_{ij}$ in
component $i$. Since $d_3^*B^U=d_3^*B^V=0$ and each relevant character
has multiplier $-2$, these actual potentials satisfy $d_1^*B^U=U$ and
$d_1^*B^V=V$. Thus the entire line $U+tV$ is coexact.

Equation~\eqref{eq:physical-dictionary} and the actual means give the same
augmented value $8+8t+k$. Scaling a violating coexact field towards zero
also excludes locally bounded quadratic penalties and the stated nonlinear
observations. Strict positivity on the ambient kernel supplies a finite
penalty for every $c<16$, which can then be restricted. The coexact field
$V$ excludes $c>16$. The same two-dimensional determinant proves the bound.
\end{proof}

\paragraph{Scope and source of the transfer.}
The pinned repository gives the actual form identities in these two sources:
\begin{quote}\small
\texttt{componentwiseFlatGradient\_eq\_hodge}\\
\path{YangMills/RG/ComponentwiseFlatHodge.lean}\\[4pt]
\texttt{actualTotalStiffness\_energy}\\
\path{YangMills/RG/ActualPacketSupportEnvelope.lean}
\end{quote}
Their proofs and the actual curl/divergence definitions were read; no new Lean compilation is
claimed. Unit 166 checks the newly constructed physical stencils rather
than relabelling a scalar diagnostic. For derivatives $nd_i$ at $n=2$,
energies and targets scale by four, so $16$ becomes $64$ with rescaled
potentials. This is finite rescaling, not a physical mass-gap claim.

The example has finite inverse moments at this one size. It refutes weak
kernel sufficiency for these prescribed rectangles at the same target;
it does not rule out another rectangular partition with the desired rank,
or establish an unbounded-size violation of the critical rank-order
conjecture. Only the trivial-background form has been identified here.


\section{Reproduction and limitations}\label{sec:evidence}
The source repository is
\href{https://github.com/lluiseriksson/THE-ERIKSSON-PROGRAMME/pull/186}{THE ERIKSSON PROGRAMME PR186},
stacked on PR185. All paths below are under
\texttt{reproducibility/matrix-moment-sweep-20261106/}.
Full source for this consolidation is pinned at
\begin{center}\small\texttt{1e3ede175efd2663f75333b3d6094fd01b0f67df}.\end{center}
Units 162 and 163 each used two distinct own CPU allocations, four actual
fresh sparse Git clones at one pinned source, and eight normal/optimized
Python processes. Outputs matched within each unit. Both runtimes were
deleted and disconnected after evidence export. Full stdout, stderr, exit
sentinels, Git logs, exact source bytes and manifests are in the corresponding
\texttt{evidence162-*.zip} and \texttt{evidence163-*.zip} archives.

\begin{center}\small
\begin{tabular}{llll}
\toprule
Unit & Executed source prefix & Fixture cases & Mutations per process\\
\midrule
162 & \texttt{f644c093} & four & seven types\\
163 & \texttt{378e879b} & four & seven types across four cases\\
\bottomrule
\end{tabular}
\end{center}
The 163 fixtures are $(D,n,h)=(3,2,1),(3,3,1),(3,3,2),(4,2,2)$.
They check actual Green and leverage identities, below/equal/above phase
inertia, equality witnesses, gap and mean-observation coercivity, and coupled
pair inertia. Unit 162 covers zero and near-zero stars, one large negative
channel and rationally rotated local eigendirections. Every acceptance check
raises explicitly on failure in both ordinary and \texttt{-O} execution.
There are 28 fixture-mutation pairs per 163 process, not 28 distinct types.
The two unit 163 recipe times were 53.83 and 29.98 seconds; requested-to-release
times were 179.71 and 102.74 seconds. Billed compute units were not exposed.

Units 165 and 166 each used two distinct own CPU allocations, four actual
fresh pinned Git clones and eight ordinary/optimized processes. Unit 165
checks the entire scalar torus matrix, the rectangular projection and three
distinct mutations per process. Unit 166 constructs the actual $72\times72$
physical matrix at $D=3,n=2,N_c=2$, with 72 unordered curl rows, 24 backward
divergence rows and 216 gradient rows. It checks the full Hodge-plus-repair
identity, both actual adjoint potentials, divergence zero and the rectangular
witnesses. Four distinct physical mutations are rejected per process: omit
divergence, double the unordered curl, reverse a potential sign and put the
field in the wrong component. Both 166 archives have 55 verified manifest
subjects. Their recipe times are 5.28 and 8.19 seconds; requested-to-release
times are 46.23 and 49.10 seconds. Both runtimes were disconnected after
export. Billed compute units remain unknown. The analytic all-$N_c$ coordinate
embedding and nonlinear limit are not extra executed cases.

Finite replay checks these finite matrices, not the universal analytical
statements. Arbitrary $q$, unequal weights, the all-size capped-family example
and the nonlinear derivative consequence are analytical arguments, not extra
executed fixtures. Earlier 159 diagnostics used source downloads rather than
fresh Git clones; their original limitation remains in their own records.
No Lean compilation certifies this appendix. A full independently prepared
audit of the imported Fourier adapter and the application form dictionary
remains necessary before broader claims.

\section{Model assessments and review context}
This appendix was prepared and checked analytically in the continuing Codex
session that prepared its sources and diagnostics. The application does not
expose the exact backend model identity for that preparation. There is no
claim of independence, blinding, human peer review or editorial acceptance.
This exact new PDF has no external model assessment and no numerical score.
An assessment of an earlier PDF does not transfer to this artifact.

The retained objections are substantive: the Cwikel mechanism and Schur
principles are classical, so independent originality is unestablished; the
observation from spectral coordinates is generally nonlocal; sparse-count
and full-moment criteria must not be conflated; a diagonal graph-rank formula
does not cover rotated eigenchannels; negative depth is neither universally
dispensable nor universally necessary; the full off-diagonal interaction must
be retained; critical kernel compatibility cannot be omitted at the target; finite diagnostics do not prove all-size statements; and no
interacting physical form equality has been supplied. These limits are part
of the result, not editorial acceptance conditions represented as satisfied.

The original positive critical rectangular construction remains open.
The common-refinement route remains stopped after three unsuccessful units.
The recorded PR29 baseline stays 20/41 with TermSource 0. This appendix adds
supporting finite mathematical results and useful counterexamples without
changing those records. It is a draft for consolidation, not a deposit,
acceptance or new verified public release.

\begin{thebibliography}{9}
\bibitem{Frank} R. L. Frank, \emph{Cwikel's theorem and the CLR inequality},
arXiv:1206.3325v1 (2012), Theorem 3.1.
\url{https://arxiv.org/abs/1206.3325v1}.
\bibitem{Zelazo} D. Zelazo and M. B\"urger, \emph{On the definiteness of the
weighted Laplacian and its connection to effective resistance},
IEEE CDC (2014), Theorem III.3.
\url{https://arxiv.org/abs/1408.2187}.
\bibitem{Spielman} D. A. Spielman, \emph{Spectral Graph Theory Lecture 14},
15 October 2018, Claim 14.6.2.
\url{https://www.cs.yale.edu/homes/spielman/561/lect14-18.pdf}.
\bibitem{Chen} W. Chen, J. Liu, Y. Chen, S. Z. Khong, D. Wang,
T. Ba\c{s}ar, L. Qiu and K. H. Johansson, \emph{Characterizing the positive
semidefiniteness of signed Laplacians via effective resistances},
IEEE CDC (2016), pp. 985--990.
\url{https://people.kth.se/~kallej/papers/network_cdc16chen.pdf}.
\bibitem{Higham} N. J. Higham, \emph{What Is the Schur Complement of a Matrix?}, 1 June 2023, Theorem 1, generalized Schur complement. \url{https://nhigham.com/2023/06/01/what-is-the-schur-complement-of-a-matrix/}.
\end{thebibliography}
\end{document}
